\documentclass[11pt]{article}
\usepackage[margin=1in]{geometry}
\usepackage{parskip}
\usepackage{hyperref}
\hypersetup{colorlinks=true, urlcolor=blue, linkcolor=black}
\sloppy

\begin{document}

\begin{center}
{\Large\bfseries Adversarial Vulnerability of Machine-Learning GNSS Spoofing Detectors to Physically Realizable Attacks}

\vspace{1em}
{\large Title Page}
\end{center}

\vspace{1.5em}
\noindent\textbf{Authors}

\begin{enumerate}
\item Joel Segun Ojerinde (Corresponding author) \\
Regional Centre for Space Science and Technology Education in Asia and the Pacific
(RCSSTEAP), Hangzhou International Innovation Institute, Beihang University,
Hangzhou 311115, China \\
Email: \href{mailto:joelojerinde@buaa.edu.cn}{joelojerinde@buaa.edu.cn}

\item Wasiu Akande Ahmed (Corresponding author) \\
Regional Centre for Space Science and Technology Education in Asia and the Pacific
(RCSSTEAP), Hangzhou International Innovation Institute, Beihang University,
Hangzhou 311115, China \\
Email: \href{mailto:ahmedwasiu24@buaa.edu.cn}{ahmedwasiu24@buaa.edu.cn}

\item Akinyinka Olukunle Akande \\
Federal University of Technology, Owerri, Nigeria \\
Email: \href{mailto:olukunle.akande@futo.edu.ng}{olukunle.akande@futo.edu.ng}
\end{enumerate}

\vspace{1em}
\noindent\textbf{Corresponding authors:} Wasiu Akande Ahmed,
\href{mailto:ahmedwasiu24@buaa.edu.cn}{ahmedwasiu24@buaa.edu.cn}; Joel Segun Ojerinde,
\href{mailto:joelojerinde@buaa.edu.cn}{joelojerinde@buaa.edu.cn}

\vspace{1.5em}
\noindent\textbf{Funding}

The authors received no specific funding for this work.

\vspace{1em}
\noindent\textbf{Authors' contributions}

J.S.O. conceived the study, developed the methodology, implemented the experiments,
analysed the results and wrote the manuscript. W.A.A. supervised the work and
reviewed the manuscript. A.O.A. reviewed the manuscript. All authors read and
approved the final manuscript.

\vspace{1em}
\noindent\textbf{Acknowledgements}

We thank the University of Texas at Austin Radionavigation Laboratory for the Texas
Spoofing Test Battery recordings and the Finnish Geospatial Research Institute for
the open-source FGI-GSRx software-defined receiver used in this study.

\end{document}
